% Náhled k~sekci o~DFS: prohledávání se noří co nejhlouběji po jedné větvi.
\begin{tikzpicture}[scale=1]

    \node[vertex, fill=maincolor!35] (V1) at (0,2)      {};
    \node[vertex, fill=maincolor!20] (V2) at (1.6,2)    {};
    \node[vertex, fill=maincolor!20] (V3) at (1.6,0.6)  {};
    \node[vertex, fill=maincolor!20] (V4) at (3.2,0.6)  {};
    \node[vertex, fill=maincolor!20] (V5) at (3.2,-0.8) {};
    \node[vertex, fill=white]        (V6) at (1.6,-0.8) {};
    \node[vertex, fill=white]        (V7) at (0,0.6)    {};
    \node[vertex, fill=white]        (V8) at (0,-0.8)   {};

    % všechny hrany grafu
    \draw[edge] (V1) -- (V7);
    \draw[edge] (V7) -- (V3);
    \draw[edge] (V3) -- (V6);
    \draw[edge] (V6) -- (V8);
    \draw[edge] (V7) -- (V8);
    \draw[edge] (V4) -- (V6);

    % větev, po které se DFS zanořuje
    \draw[edge, very thick, darkred] (V1) -- (V2);
    \draw[edge, very thick, darkred] (V2) -- (V3);
    \draw[edge, very thick, darkred] (V3) -- (V4);
    \draw[edge, very thick, darkred] (V4) -- (V5);

    \draw[->, >=stealth, thick, darkred, dashed]
        (4.0,2) to[bend left=45] node[midway, right] {\small hloubka} (4.0,-0.8);

\end{tikzpicture}
